\documentclass[a4paper,12pt,twoside]{memoir}

% Castellano
\usepackage[spanish,es-tabla]{babel}
\selectlanguage{spanish}
\usepackage[utf8]{inputenc}
\usepackage[T1]{fontenc}
\usepackage{lmodern} % Scalable font
\usepackage{microtype}
\usepackage{placeins}

\RequirePackage{booktabs}
\RequirePackage[table]{xcolor}
\RequirePackage{xtab}
\RequirePackage{multirow}

% Links
\PassOptionsToPackage{hyphens}{url}\usepackage[colorlinks]{hyperref}
\hypersetup{
	allcolors = {red}
}

% Ecuaciones
\usepackage{amsmath}

% Rutas de fichero / paquete
\newcommand{\ruta}[1]{{\sffamily #1}}

% Párrafos
\nonzeroparskip

% Huérfanas y viudas
\widowpenalty100000
\clubpenalty100000

% Imagenes
\usepackage{graphicx}
\newcommand{\imagen}[2]{
	\begin{figure}[!h]
		\centering
		\includegraphics[width=0.9\textwidth]{#1}
		\caption{#2}\label{fig:#1}
	\end{figure}
	\FloatBarrier
}

\newcommand{\imagenflotante}[2]{
	\begin{figure}%[!h]
		\centering
		\includegraphics[width=0.9\textwidth]{#1}
		\caption{#2}\label{fig:#1}
	\end{figure}
}



% El comando \figura nos permite insertar figuras comodamente, y utilizando
% siempre el mismo formato. Los parametros son:
% 1 -> Porcentaje del ancho de página que ocupará la figura (de 0 a 1)
% 2 --> Fichero de la imagen
% 3 --> Texto a pie de imagen
% 4 --> Etiqueta (label) para referencias
% 5 --> Opciones que queramos pasarle al \includegraphics
% 6 --> Opciones de posicionamiento a pasarle a \begin{figure}
\newcommand{\figuraConPosicion}[6]{%
  \setlength{\anchoFloat}{#1\textwidth}%
  \addtolength{\anchoFloat}{-4\fboxsep}%
  \setlength{\anchoFigura}{\anchoFloat}%
  \begin{figure}[#6]
    \begin{center}%
      \Ovalbox{%
        \begin{minipage}{\anchoFloat}%
          \begin{center}%
            \includegraphics[width=\anchoFigura,#5]{#2}%
            \caption{#3}%
            \label{#4}%
          \end{center}%
        \end{minipage}
      }%
    \end{center}%
  \end{figure}%
}

%
% Comando para incluir imágenes en formato apaisado (sin marco).
\newcommand{\figuraApaisadaSinMarco}[5]{%
  \begin{figure}%
    \begin{center}%
    \includegraphics[angle=90,height=#1\textheight,#5]{#2}%
    \caption{#3}%
    \label{#4}%
    \end{center}%
  \end{figure}%
}
% Para las tablas
\newcommand{\otoprule}{\midrule [\heavyrulewidth]}
%
% Nuevo comando para tablas pequeñas (menos de una página).
\newcommand{\tablaSmall}[5]{%
 \begin{table}
  \begin{center}
   \rowcolors {2}{gray!35}{}
   \begin{tabular}{#2}
    \toprule
    #4
    \otoprule
    #5
    \bottomrule
   \end{tabular}
   \caption{#1}
   \label{tabla:#3}
  \end{center}
 \end{table}
}

%
% Nuevo comando para tablas pequeñas (menos de una página).
\newcommand{\tablaSmallSinColores}[5]{%
 \begin{table}[H]
  \begin{center}
   \begin{tabular}{#2}
    \toprule
    #4
    \otoprule
    #5
    \bottomrule
   \end{tabular}
   \caption{#1}
   \label{tabla:#3}
  \end{center}
 \end{table}
}

\newcommand{\tablaApaisadaSmall}[5]{%
\begin{landscape}
  \begin{table}
   \begin{center}
    \rowcolors {2}{gray!35}{}
    \begin{tabular}{#2}
     \toprule
     #4
     \otoprule
     #5
     \bottomrule
    \end{tabular}
    \caption{#1}
    \label{tabla:#3}
   \end{center}
  \end{table}
\end{landscape}
}

%
% Nuevo comando para tablas grandes con cabecera y filas alternas coloreadas en gris.
\newcommand{\tabla}[6]{%
  \begin{center}
    \tablefirsthead{
      \toprule
      #5
      \otoprule
    }
    \tablehead{
      \multicolumn{#3}{l}{\small\sl continúa desde la página anterior}\\
      \toprule
      #5
      \otoprule
    }
    \tabletail{
      \hline
      \multicolumn{#3}{r}{\small\sl continúa en la página siguiente}\\
    }
    \tablelasttail{
      \hline
    }
    \bottomcaption{#1}
    \rowcolors {2}{gray!35}{}
    \begin{xtabular}{#2}
      #6
      \bottomrule
    \end{xtabular}
    \label{tabla:#4}
  \end{center}
}

%
% Nuevo comando para tablas grandes con cabecera.
\newcommand{\tablaSinColores}[6]{%
  \begin{center}
    \tablefirsthead{
      \toprule
      #5
      \otoprule
    }
    \tablehead{
      \multicolumn{#3}{l}{\small\sl continúa desde la página anterior}\\
      \toprule
      #5
      \otoprule
    }
    \tabletail{
      \hline
      \multicolumn{#3}{r}{\small\sl continúa en la página siguiente}\\
    }
    \tablelasttail{
      \hline
    }
    \bottomcaption{#1}
    \begin{xtabular}{#2}
      #6
      \bottomrule
    \end{xtabular}
    \label{tabla:#4}
  \end{center}
}

%
% Nuevo comando para tablas grandes sin cabecera.
\newcommand{\tablaSinCabecera}[5]{%
  \begin{center}
    \tablefirsthead{
      \toprule
    }
    \tablehead{
      \multicolumn{#3}{l}{\small\sl continúa desde la página anterior}\\
      \hline
    }
    \tabletail{
      \hline
      \multicolumn{#3}{r}{\small\sl continúa en la página siguiente}\\
    }
    \tablelasttail{
      \hline
    }
    \bottomcaption{#1}
  \begin{xtabular}{#2}
    #5
   \bottomrule
  \end{xtabular}
  \label{tabla:#4}
  \end{center}
}



\definecolor{cgoLight}{HTML}{EEEEEE}
\definecolor{cgoExtralight}{HTML}{FFFFFF}

%
% Nuevo comando para tablas grandes sin cabecera.
\newcommand{\tablaSinCabeceraConBandas}[5]{%
  \begin{center}
    \tablefirsthead{
      \toprule
    }
    \tablehead{
      \multicolumn{#3}{l}{\small\sl continúa desde la página anterior}\\
      \hline
    }
    \tabletail{
      \hline
      \multicolumn{#3}{r}{\small\sl continúa en la página siguiente}\\
    }
    \tablelasttail{
      \hline
    }
    \bottomcaption{#1}
    \rowcolors[]{1}{cgoExtralight}{cgoLight}

  \begin{xtabular}{#2}
    #5
   \bottomrule
  \end{xtabular}
  \label{tabla:#4}
  \end{center}
}



\graphicspath{ {./img/} }

% Capítulos
\chapterstyle{bianchi}
\newcommand{\capitulo}[2]{
	\setcounter{chapter}{#1}
	\setcounter{section}{0}
	\setcounter{figure}{0}
	\setcounter{table}{0}
	\chapter*{#2}
	\addcontentsline{toc}{chapter}{#2}
	\markboth{#2}{#2}
}

% Apéndices
\renewcommand{\appendixname}{Apéndice}
\renewcommand*\cftappendixname{\appendixname}

\newcommand{\apendice}[1]{
	%\renewcommand{\thechapter}{A}
	\chapter{#1}
}

\renewcommand*\cftappendixname{\appendixname\ }

% Formato de portada
\makeatletter
\usepackage{xcolor}
\newcommand{\tutor}[1]{\def\@tutor{#1}}
\newcommand{\course}[1]{\def\@course{#1}}
\definecolor{cpardoBox}{HTML}{E6E6FF}
\def\maketitle{
  \null
  \thispagestyle{empty}
  % Cabecera ----------------
\noindent\includegraphics[width=\textwidth]{cabecera}\vspace{1cm}%
  \vfill
  % Título proyecto y escudo informática ----------------
  \colorbox{cpardoBox}{%
    \begin{minipage}{.8\textwidth}
      \vspace{.5cm}\Large
      \begin{center}
      \textbf{TFG del Grado en Ingeniería Informática}\vspace{.6cm}\\
      \textbf{\LARGE\@title{}}
      \end{center}
      \vspace{.2cm}
    \end{minipage}

  }%
  \hfill\begin{minipage}{.20\textwidth}
    \includegraphics[width=\textwidth]{escudoInfor}
  \end{minipage}
  \vfill
  % Datos de alumno, curso y tutores ------------------
  \begin{center}%
  {%
    \noindent\LARGE
    Presentado por \@author{}\\ 
    en Universidad de Burgos --- \@date{}\\
    Tutores: \@tutor{}\\
  }%
  \end{center}%
  \null
  \cleardoublepage
  }
\makeatother

\newcommand{\nombre}{Ismael Franco Hernando} %%% cambio de comando

% Datos de portada
\title{X-Teeth}
\author{\nombre}
\tutor{César Ignacio García Osorio y \\José Miguel Ramírez Sanz}
\date{\today}

\begin{document}

\maketitle


\newpage\null\thispagestyle{empty}\newpage


%%%%%%%%%%%%%%%%%%%%%%%%%%%%%%%%%%%%%%%%%%%%%%%%%%%%%%%%%%%%%%%%%%%%%%%%%%%%%%%%%%%%%%%%
\thispagestyle{empty}


\noindent\includegraphics[width=\textwidth]{cabecera}\vspace{1cm}

\noindent D. César Ignacio García Osorio, profesor del departamento de  Ingeniería Informática, área de  Lenguajes y Sistemas Informáticos y D. José Miguel Ramírez Sanz, profesor del departamento de Ingenería Informática, área de Lenguajes y Sistemas Informáticos.

\noindent Exponen:

\noindent Que el alumno D. \nombre, con DNI 71363577M, ha realizado el Trabajo final de Grado en Ingeniería Informática titulado X-Teeth. 

\noindent Y que dicho trabajo ha sido realizado por el alumno bajo la dirección del que suscribe, en virtud de lo cual se autoriza su presentación y defensa.

\begin{center} %\large
En Burgos, {\large \today}
\end{center}

\vfill\vfill\vfill

% Author and supervisor
\begin{minipage}{0.45\textwidth}
\begin{flushleft} %\large
Vº. Bº. del Tutor:\\[2cm]
D. César Ignacio García Osorio
\end{flushleft}
\end{minipage}
\hfill
\begin{minipage}{0.45\textwidth}
\begin{flushleft} %\large
Vº. Bº. del co-tutor:\\[2cm]
D. José Miguel Ramírez Sanz
\end{flushleft}
\end{minipage}
\hfill

\vfill

% para casos con solo un tutor comentar lo anterior
% y descomentar lo siguiente
%Vº. Bº. del Tutor:\\[2cm]
%D. nombre tutor


\newpage\null\thispagestyle{empty}\newpage




\frontmatter

% Abstract en castellano
\renewcommand*\abstractname{Resumen}
\begin{abstract}
En este primer apartado se hace una \textbf{breve} presentación del tema que se aborda en el proyecto.
\end{abstract}

\renewcommand*\abstractname{Descriptores}
\begin{abstract}
Palabras separadas por comas que identifiquen el contenido del proyecto Ej: servidor web, buscador de vuelos, android \ldots
\end{abstract}

\clearpage

% Abstract en inglés
\renewcommand*\abstractname{Abstract}
\begin{abstract}
A \textbf{brief} presentation of the topic addressed in the project.
\end{abstract}

\renewcommand*\abstractname{Keywords}
\begin{abstract}
keywords separated by commas.
\end{abstract}

\clearpage

% Indices
\tableofcontents

\clearpage

\listoffigures

\clearpage

\listoftables
\clearpage

\mainmatter
\capitulo{1}{Introducción}

Descripción del contenido del trabajo y del estrucutra de la memoria y del resto de materiales entregados.

\capitulo{2}{Objetivos del proyecto}

Este apartado contiene los objetivos del proyecto. Se pueden distinguir dos tipos: generales y técnicos. Dichos objetivos buscan cumplirse con el desarrollo del propio proyecto.

\section{Objetivos Generales}
Los objetivos generales del proyecto son:
\begin{itemize}
    \item Investigar y aplicar nuevas funcionalidades del \emph{deep learning} en la rama de la odontología.
    \item Desarrollar una aplicación web, sencilla de usar y con un fácil acceso para todo el mundo.
    \item Ayudar en el proceso de la endodoncia, de forma que la obtención de la longitud, actualmente hecha de forma manual, se pueda automatizar con la aplicación.
    \item Intentar conseguir que el error de la aplicación en la estimación de la longitud del diente no sea superior a 0.5 milímetros.
\end{itemize}

\section{Objetivos Técnicos}
Los objetivos técnicos del proyecto son:
\begin{itemize}
    \item Crear una aplicación web capaz de calcular la longitud de un diente a través de una radiografía.
    \item Permitir el uso de la aplicación de dos modos distintos:
    \begin{itemize}
        \item Acceder directamente al \emph{notebook} en \emph{Google Colab}.
        \item Usar \emph{ngrok} para acceder a la dirección web que contiene el \emph{notebook} con la aplicación en Gamma.
    \end{itemize}
    \item Usar el lenguaje de programación \emph{Python} para realizar las pruebas pertinentes y la aplicación final, junto con las distintas librerías que existen en \emph{Python}.
    \item Conseguir que los tiempos de espera en la aplicación sean los más breves posibles.
    \item Utilizar la plataforma de \emph{GitHub}, de forma que se vea reflejado el trabajo incremental que se lleva a cabo en cada \emph{sprint}.
\end{itemize}

\section{Objetivos Personales}
Los objetivos marcados a nivel personal son:
\begin{itemize}
    \item Aprender sobre el proceso de la endodoncia, para poder investigar y desarrollar una aplicación que permita ayudar a los odontólogos en su trabajo.
    \item Aplicar los conocimientos adquiridos a lo largo del grado universitario.
    \item Mejorar los conocimientos de \emph{Python} y aplicar aquellos ya adquiridos.
    \item Aprender sobre el \emph{deep learning} para conocer sus posibilidades y poder aplicarlas a lo largo del proyecto.
\end{itemize}
\capitulo{3}{Conceptos teóricos}

En aquellos proyectos que necesiten para su comprensión y desarrollo de unos conceptos teóricos de una determinada materia o de un determinado dominio de conocimiento, debe existir un apartado que sintetice dichos conceptos.

Algunos conceptos teóricos de \LaTeX \footnote{Créditos a los proyectos de Álvaro López Cantero: Configurador de Presupuestos y Roberto Izquierdo Amo: PLQuiz}.

\section{Secciones}

Las secciones se incluyen con el comando section.

\subsection{Subsecciones}

Además de secciones tenemos subsecciones.

\subsubsection{Subsubsecciones}

Y subsecciones. 


\section{Referencias}

Las referencias se incluyen en el texto usando cite \cite{wiki:latex}. Para citar webs, artículos o libros \cite{koza92}.


\section{Imágenes}

Se pueden incluir imágenes con los comandos standard de \LaTeX, pero esta plantilla dispone de comandos propios como por ejemplo el siguiente:

\imagen{escudoInfor}{Autómata para una expresión vacía}



\section{Listas de items}

Existen tres posibilidades:

\begin{itemize}
	\item primer item.
	\item segundo item.
\end{itemize}

\begin{enumerate}
	\item primer item.
	\item segundo item.
\end{enumerate}

\begin{description}
	\item[Primer item] más información sobre el primer item.
	\item[Segundo item] más información sobre el segundo item.
\end{description}
	
\begin{itemize}
\item 
\end{itemize}

\section{Tablas}

Igualmente se pueden usar los comandos específicos de \LaTeX o bien usar alguno de los comandos de la plantilla.

\tablaSmall{Herramientas y tecnologías utilizadas en cada parte del proyecto}{l c c c c}{herramientasportipodeuso}
{ \multicolumn{1}{l}{Herramientas} & App AngularJS & API REST & BD & Memoria \\}{ 
HTML5 & X & & &\\
CSS3 & X & & &\\
BOOTSTRAP & X & & &\\
JavaScript & X & & &\\
AngularJS & X & & &\\
Bower & X & & &\\
PHP & & X & &\\
Karma + Jasmine & X & & &\\
Slim framework & & X & &\\
Idiorm & & X & &\\
Composer & & X & &\\
JSON & X & X & &\\
PhpStorm & X & X & &\\
MySQL & & & X &\\
PhpMyAdmin & & & X &\\
Git + BitBucket & X & X & X & X\\
Mik\TeX{} & & & & X\\
\TeX{}Maker & & & & X\\
Astah & & & & X\\
Balsamiq Mockups & X & & &\\
VersionOne & X & X & X & X\\
} 

\capitulo{4}{Técnicas y herramientas}

Esta parte de la memoria tiene como objetivo presentar las técnicas metodológicas y las herramientas de desarrollo que se han utilizado para llevar a cabo el proyecto. Si se han estudiado diferentes alternativas de metodologías, herramientas, bibliotecas se puede hacer un resumen de los aspectos más destacados de cada alternativa, incluyendo comparativas entre las distintas opciones y una justificación de las elecciones realizadas. 
No se pretende que este apartado se convierta en un capítulo de un libro dedicado a cada una de las alternativas, sino comentar los aspectos más destacados de cada opción, con un repaso somero a los fundamentos esenciales y referencias bibliográficas para que el lector pueda ampliar su conocimiento sobre el tema.



\capitulo{5}{Aspectos relevantes del desarrollo del proyecto}

Este apartado pretende recoger los aspectos más interesantes del desarrollo del proyecto, comentados por los autores del mismo.
Debe incluir desde la exposición del ciclo de vida utilizado, hasta los detalles de mayor relevancia de las fases de análisis, diseño e implementación.
Se busca que no sea una mera operación de copiar y pegar diagramas y extractos del código fuente, sino que realmente se justifiquen los caminos de solución que se han tomado, especialmente aquellos que no sean triviales.
Puede ser el lugar más adecuado para documentar los aspectos más interesantes del diseño y de la implementación, con un mayor hincapié en aspectos tales como el tipo de arquitectura elegido, los índices de las tablas de la base de datos, normalización y desnormalización, distribución en ficheros3, reglas de negocio dentro de las bases de datos (EDVHV GH GDWRV DFWLYDV), aspectos de desarrollo relacionados con el WWW...
Este apartado, debe convertirse en el resumen de la experiencia práctica del proyecto, y por sí mismo justifica que la memoria se convierta en un documento útil, fuente de referencia para los autores, los tutores y futuros alumnos.

\capitulo{6}{Trabajos relacionados}

\section{Application of Deep Learning in Dentistry and
Implantology}
Autores: Dae-Young Kang, Hieu Pham Duong y Jung-Chul Park

Este trabajo \cite{kang2020application} recoge los resultados de la investigación de algoritmos existentes de \emph{deep learning} en la rama de la odontología. Es una colaboración entre las Universidades de Dankook y de Vietnam, publicado en el año 2020.

Lo primero que resaltan es que en la odontología se requiere una precisión muy elevada, ya que se trabaja con medidas milimétricas y un error humano puede tener consecuencias irreparables. Es por ello, que defienden que la inteligencia artificial y el \emph{deep learning} suponen una gran mejora de precisión y de detección de anomalías que pueden ser indetectables a la vista humana.

En la parte de investigación se habla sobre la diversidad de aplicaciones que tiene el \emph{deep learning} en la odontología, y una de ellas es la segmentación e identificación de dientes (tarea que se ha llevado a cabo en ese proyecto), aunque no dice nada de cálculo de longitud de dientes.

La conclusión final del artículo es que se espera que la utilización de algoritmos de \emph{deep learning} mejore los tiempos de trabajo de los dentistas y los resultados de los diferentes tratamientos, ayudando a los dentistas en casi todos los aspectos prácticos, como detección de dientes, detección y clasificación de enfermedades, evaluación de prótesis y reducción de errores de precisión.

\section{Software Second-Opinion}
La aplicación \emph{Second-Opinion}\footnote{Second-Opinion: \url{https://www.hellopearl.com/products/second-opinion}} es un \emph{software} de pago que permite a los odontólogos que la utilizan cargar sus radiografías y la aplicación, a través de la inteligencia artificial, calcula o muestra los valores que los odontólogos necesitan.

Las principales tareas que realiza la aplicación son:
\begin{itemize}
    \item Detección de enfermedades.
    \item Validación de implantes.
    \item Segmentación de dientes.
    \item Obtención del canal radicular.
\end{itemize}

Pese a ser un software muy completo, no incluye el cálculo de la longitud del diente, cosa que sí que realiza la aplicación de este proyecto.

Para finalizar, este \emph{software} ha sido aprobado por la FDA (\emph{Food and Drugs Administration}, Administración de Alimentos y Medicamentos en español, en Estados Unidos)\footnote{Artículo aprobación software: \url{https://www.hellopearl.com/press-release/fda-clears-worlds-first-ai-software-to-read-dental-x-rays}}, para que los dentistas del país puedan utilizarlo con los pacientes. Además, esta aplicación ha sido la primera de este tipo en aprobarse en Estados Unidos por lo que se puede ver como es una rama en la que aún queda mucho por explorar.
\capitulo{7}{Conclusiones y Líneas de trabajo futuras}

\section{Conclusiones}
Este proyecto surge de la mano del odontólogo Álvaro Zubizarreta Macho, quien propuso la creación de una aplicación que le ayudara en la endodoncia al ser un proceso totalmente artesanal y con una gran posibilidad a fallo.

El primer problema que hubo fue crear un modelo con tan pocas radiografías, pero, pese a ello, se solucionó transformando las radiografías prestadas, gracias al \emph{data augmentation}.

Tras conseguir el modelo se mostraron los primeros resultados de la segmentación de dientes y nervios al odontólogo, quien reconoció que eran bastantes buenos. Este profesional fue de ayuda en lo relacionado con las técnicas de cálculo, permitiendo encontrar la forma más óptima. Lamentablemente, los resultados obtenidos no se han podido probar en su totalidad debido a diversos problemas en la resolución y tamaño de las radiografías, los cuales no han sido solucionados a lo largo del proyecto.

A pesar de esto, los resultados son gratamente satisfactorios debido a que el odontólogo afirmó que las predicciones del modelo eran correctas. Sin embargo, hubiera sido más gratificante poder analizar la calidad de los resultados.
Durante este proceso se ha disfrutado del aprendizaje del \emph{deep learning} en \emph{Python}, debido a ser algo de lo que no se tenía mucho conocimiento. Además, se han conocido una gran cantidad de herramientas útiles de cara al futuro, como \emph{Ngrok} y \emph{tmux}.

Finalmente, cabe destacar que se ha desarrollado una aplicación sencilla y fácil de usar para usuarios alejados del mundo de la programación y los \emph{notebooks}, permitiendo que esta herramienta pueda ser usada por cualquier cliente.

\section{Líneas de trabajo futuras}
En mi opinión, unas mejoras que se podrían realizar al proyecto en el futuro serían:

\begin{enumerate}
    \item \textbf{Cálculo de longitud de los demás dientes de la radiografía:} actualmente se devuelve la longitud del diente y del nervio con mayor precisión por parte del modelo de \emph{Detectron2}. Es por ello, que se podría calcular la longitud de todos los dientes que aparezcan en las radiografías con el fin de poder ayudar a los odontólogos en caso de que necesiten más longitudes.
    
    \item \textbf{Aumentar el número de radiografías con el que entrenar al modelo:} cómo se ha comentado en puntos anteriores, el proyecto tan solo ha contado con 10 radiografías, las cuales han tenido que ser transformadas para poder contar con un gran número de imágenes con las que entrenar al modelo. Pero realmente, no dejan de ser 10 imágenes reales con sus variaciones. Es por ello, que sería buena idea aumentar considerablemente el número de radiografías para conseguir un modelo más robusto y capaz de detectar mejor los dientes y nervios en aquellos casos más extraños.
    
    \item \textbf{Desarrollar una aplicación web:} la aplicación final se ha desarrollado en un \emph{notebook}, por lo que sería buena idea transportarla a una aplicación web de forma que resulte más cómodo el acceso a los clientes. Evitando, por ejemplo, que tengan que esperar tiempos de espera en la instalación de \emph{Detectron2} en \emph{Google Colab}.
    
    \item \textbf{Añadir nuevas funcionalidades a la aplicación:} el proyecto se ha centrado únicamente en crear una aplicación capaz de calcular la longitud del diente, pero en el futuro se podrían implementar más funcionalidades con las que ayudar en las endodoncias o incluso en alguna otra rama de la odontología.
\end{enumerate}



\bibliographystyle{plain}
\bibliography{bibliografia}

\end{document}
